\documentclass[12pt,a4paper]{article}
\usepackage[utf8]{inputenc}
\usepackage[indonesian]{babel}
\usepackage{amsmath}
\usepackage{amsfonts}
\usepackage{amssymb}
\usepackage{geometry}

\geometry{margin=2.5cm}

\title{Menghitung Massa dan Pusat Massa Kawat Linier}
\author{Ahmad Fakhrul Bawani}
\date{}

\begin{document}

\maketitle

\section*{Soal}

Find the mass $M$ and center of mass $\bar{x}$ of the linear wire covering the given interval and having the given density $\delta(x)$:
\begin{equation*}
  -2 \le x \le 2, \quad \delta(x) = 4 + 5x^4
\end{equation*}

\subsection*{1. Rumus Umum}
Untuk kawat linier yang terletak pada interval $[a, b]$ dengan fungsi kerapatan $\delta(x)$:
\begin{itemize}
  \item \textbf{Massa Total ($M$):}
    \begin{equation*}
      M = \int_{a}^{b} \delta(x) \, dx
    \end{equation*}
  \item \textbf{Momen Pertama terhadap Titik Asal ($M_0$):}
    \begin{equation*}
      M_0 = \int_{a}^{b} x \cdot \delta(x) \, dx
    \end{equation*}
  \item \textbf{Pusat Massa ($\bar{x}$):}
    \begin{equation*}
      \bar{x} = \frac{M_0}{M}
    \end{equation*}
\end{itemize}

\subsection*{2. Solusi Menghitung Massa ($M$)}
Substitusikan fungsi kerapatan ke dalam rumus massa dengan batas $a = -2$ dan $b = 2$:
\begin{align*}
  M &= \int_{-2}^{2} (4 + 5x^4) \, dx
\end{align*}
Karena fungsi $\delta(x) = 4 + 5x^4$ merupakan fungsi genap ($\delta(-x) = \delta(x)$) dan diintegrasikan pada batas yang simetris, kita dapat menyederhanakannya menjadi:
\begin{align*}
  M &= 2 \int_{0}^{2} (4 + 5x^4) \, dx \\
  M &= 2 \left[ 4x + x^5 \right]_{0}^{2} \\
  M &= 2 \left( \left[ 4(2) + 2^5 \right] - 0 \right) \\
  M &= 2 (8 + 32) \\
  M &= 2 (40) \\
  M &= 80
\end{align*}

\subsection*{3. Solusi Menghitung Pusat Massa ($\bar{x}$)}
Pertama, hitung momen pertama $M_0$:
\begin{align*}
  M_0 &= \int_{-2}^{2} x(4 + 5x^4) \, dx \\
  M_0 &= \int_{-2}^{2} (4x + 5x^5) \, dx
\end{align*}
Perhatikan bahwa fungsi $g(x) = 4x + 5x^5$ adalah fungsi ganjil ($g(-x) = -g(x)$). Berdasarkan sifat integral tentu, mengintegrasikan fungsi ganjil pada batas simetris $[-a, a]$ akan selalu menghasilkan nilai nol:
\begin{equation*}
  M_0 = 0
\end{equation*}

Sekarang, kita tentukan titik pusat massanya ($\bar{x}$):
\begin{equation*}
  \bar{x} = \frac{M_0}{M} = \frac{0}{80} = 0
\end{equation*}

\subsection*{Kesimpulan}
Jadi, hasil perhitungan untuk kawat linier tersebut adalah:
\begin{itemize}
  \item Massa total kawat ($M$) = \textbf{$80$}
  \item Pusat massa kawat ($\bar{x}$) = \textbf{$0$}
\end{itemize}

\end{document}